\section{Discussion}
\label{sec:discussion}

\subsection{The Aggregate-Security Trade-off}
\label{sec:discussion_tradeoff}

The central empirical finding is a trade-off between aggregate F1-macro and
security-relevant recall on rare stealthy classes. The hybrid OCSVM~+~XGBoost
pipeline achieves backdoor recall $= 0.925$ versus $0.362$ for standalone
XGBoost ($+56.4\%$), at the cost of $-13.4\%$ F1-macro.

This trade-off is operationally rational in the ICS context: a missed backdoor
event (false negative) may enable persistent adversarial access preceding a
destructive attack sequence, while a false alarm merely triggers an analyst
review. Under this asymmetric cost structure, recall on backdoor and ransomware
classes constitutes the primary design criterion, and F1-macro serves as a
secondary constraint.

\subsection{Why ADASYN Degrades Overall Performance (A2)}
\label{sec:discussion_adasyn}

The ablation result for A2 (ADASYN without filtering, $-8.1\%$ F1-macro) is
counterintuitive. ADASYN generates synthetic samples to equalize class
distribution: on TON\,IoT with 10 classes and scanning/ddos dominating,
ADASYN inflates the training set from $321{,}631$ to $599{,}072$ samples
($+86\%$). The synthetic interpolations in the feature space of mitm
(212 samples) and ransomware (273 samples) are geometrically unreliable
given the small neighborhood sizes, and may introduce mislabelled boundary
regions that confound XGBoost's splits. This finding motivates the use of
cost-sensitive learning (class weights) rather than synthetic oversampling
for intrusion detection datasets.

\subsection{LightGBM Underperformance: A Methodological Warning}
\label{sec:discussion_lgb}

LightGBM achieves F1-macro$=0.4141$ despite \texttt{class\_weight='balanced'},
compared to $0.8039$ in the 10-seed evaluation (subsample). Investigation
reveals that LightGBM silently ignores the \texttt{class\_weight} parameter
when input data are \texttt{numpy} arrays rather than \texttt{pandas}
DataFrames in the evaluated version. This constitutes a critical reproducibility
pitfall: the same configuration produces qualitatively different results
depending on input format.

\subsection{Limitations and Threats to Validity}
\label{sec:discussion_limits}

\textbf{Single dataset.} All results are obtained on one stratified sample
from TON\,IoT. Validation on SWaT~\cite{ref_SWaT} or BATADAL~\cite{ref_BATADAL}
is required to assess generalization to pure ICS environments.

\textbf{Elevated attack proportion.} The sampled TON\,IoT subset (85\% attacks)
differs structurally from operational ICS networks where normal traffic often
exceeds 90\%. This limits OCSVM's discriminating power and enrichment factor.

\textbf{H3 untested.} Deep learning models (CNN-LSTM, PatchTST) were not
evaluated due to computational constraints. These architectures are expected
to improve recall on temporally-structured attacks (Replay, Modify Parameter).

\textbf{Static threshold $\theta$.} OCSVM's threshold is fixed after validation
calibration and does not adapt to concept drift in production ICS traffic.
